\documentclass[10pt,a4paper]{article}
\newcommand{\CommonPath}{../..}
\usepackage[english,noicons]{\CommonPath/scholia/scholia}

% Test: noicons option — icon helper degrades cleanly

\begin{document}

\section*{Noicons Mode Tests}

% --- scholiaicon with fallback ---
Icon with fallback: \scholiaicon[fallback={[bulb]}]{lightbulb.pdf}

% --- scholiaicon without fallback (should produce nothing) ---
Icon without fallback: [\scholiaicon{lightbulb.pdf}] (brackets should be adjacent)

% --- cardbox should render without icon badge ---
\begin{cardbox}[title={Card without icons}, icon={lightbulb.pdf}, color={blue}]
  The badge should not render in noicons mode.
  The floating title label should still appear.
\end{cardbox}

% --- existing callout family should also stay icon-free ---
\begin{tipbox}[Tip without icon]
  Existing callouts should not load icon assets in noicons mode.
\end{tipbox}

% --- exercise titles should render without icon prefixes ---
\begin{exguided}[Guided exercise without icon]
  Exercise titles should not reserve icon space in noicons mode.
\end{exguided}

% --- ideacard preset ---
\begin{ideacard}[title={Idea without badge}]
  The idea card should render its body and title but skip the badge icon.
\end{ideacard}

% --- typed page banner ---
\newTypedPage{chapter-1}{notion}{Functions}{Images and antecedents}
Typed-page banners should render without icon circles in noicons mode.

% --- summary references and index banners ---
Summary reference without icon badge: \summaryRef{concepts}{N1}[Definition of a function]

\begin{indexTypeTable}*{concepts}
  \indexEntry{N1}{Definition without icon}
\end{indexTypeTable}

\end{document}
